Fix \(\chi>1\). Along the triangular center sequence, every observational cell of \((A,X,Z_{M_d})\) has probability bounded below by a positive constant depending only on \(\chi\), uniformly in \(m\). Indeed, the observational treatment kernels lie in \([3/8,5/8]\), the early signal belongs to \((0,10^{-3})\), the late signal \(y_\chi\) is fixed strictly below one, and the remaining signals equal \(1/2\). Also \(P(G=O)=1/2\). Standard binomial concentration applied to the finitely many cells therefore shows that every training and validation cell count is of order \(m\) with probability tending to one, and the probability of any zero denominator is exponentially small.

For a cell \(v=(a,x,z_{M_d})\), write
\[
 f_{d,m}^*(v)=\mathbb E_{P_{\chi,\epsilon_m}}(Y\mid A,X,Z_{M_d}=v,G=O).
\]
Both \(Y\) and every saturated cell mean lie in \([-1,1]\). Conditional on the cell counts, each nonempty-cell mean is an average of bounded independent variables, and its count is of order \(m\). A finite union bound over both designs and their finitely many cells gives
\[
 \max_{d\in\mathcal D}\max_v
 |\widehat f_{d,m}^A(v)-f_{d,m}^*(v)|=O_P(m^{-1/2}).        \tag{43}
\]
Conditional on the training half, the validation losses are independent and bounded by four. The same finite-class concentration argument gives
\[
 \max_d\left|\widehat R_{d,m}^A-
 \mathbb E\left[
 \{Y-\widehat f_{d,m}^A(A,X,Z_{M_d})\}^2
 \mathrel{\Big|}G=O,\mathcal I_{m,O}^{\mathrm{tr}}
 \right]\right|=O_P(m^{-1/2}).                            \tag{44}
\]
Because \(f_{d,m}^*\) is the conditional-mean projection, the cross term vanishes cell by cell and hence
\[
\begin{aligned}
 &\mathbb E\left[
 \{Y-\widehat f_{d,m}^A(A,X,Z_{M_d})\}^2
 \mathrel{\Big|}G=O,\mathcal I_{m,O}^{\mathrm{tr}}
 \right]-R_d^A(P_{\chi,\epsilon_m})\\
 &\qquad=
 \mathbb E\left[
 \{\widehat f_{d,m}^A(A,X,Z_{M_d})-f_{d,m}^*(A,X,Z_{M_d})\}^2
 \mathrel{\Big|}G=O,\mathcal I_{m,O}^{\mathrm{tr}}
 \right]
 =O_P(m^{-1}).                                             \tag{45}
\end{aligned}
\]
Here the last equality follows from (43) and the finite number of cells. Combining (44)--(45) proves (1).

Lemma~\(\mathrm{lem{:}treatment\text{-}aware\text{-}sharp\text{-}binary\text{-}elbow}\) gives, uniformly along the triangular sequence,
\[
 R_{d_{\mathrm p}}^A(P_{\chi,\epsilon_m})<\frac{3}{10},
 \qquad
 R_{d_{\mathrm e}}^A(P_{\chi,\epsilon_m})=\frac{585}{1024}.
\]
The population gap is bounded below by \(585/1024-3/10>0\). The uniform empirical error in (1), together with deterministic tie breaking, therefore implies
\[
 P_{\chi,\epsilon_m}\{\widehat d_m^A=d_{\mathrm p}\}\longrightarrow1,
\]
which is (2).

On this event, Lemma~\(\mathrm{lem{:}sharp\text{-}uniform\text{-}binary\text{-}elbow}\) identifies \(d_{\mathrm e}\) as the unique minimizer of both profiles and gives
\[
 \frac{\mathcal V_{d_{\mathrm p}}^*(P_{\chi,\epsilon_m})}
 {\mathcal V_{d_{\mathrm e}}^*(P_{\chi,\epsilon_m})}
 =\frac{\mathcal V_{d_{\mathrm p},\mathrm{IKM}}^*(P_{\chi,\epsilon_m})}
 {\mathcal V_{d_{\mathrm e},\mathrm{IKM}}^*(P_{\chi,\epsilon_m})}
 >\frac{\epsilon_m^{-2}+14}{726}\longrightarrow\infty.  \tag{46}
\]
Since the event has probability tending to one, both random achieved ratios diverge in probability, proving (3). Assumption~\(\mathrm{ass{:}independent\text{-}pilot}\) makes this pilot-only rule feasible before the main-study measurements are chosen. Its construction uses bounded finite-cell averages and no bridge inverse, inverse-problem learner, or growing sieve.